% GENERATED FILE. Edit canonical lessons, then run npm run generate:books.
% canonical-source-hash: fnv1a64:ab94f207d7c9c56b
% canonical-lessons: KA-C489-ibbani, KA-C489-kulirgali, KA-C489-manjugadde, KA-C489-heppugattu, KA-C489-minugu, KA-R489-review-483-489, KA-R489-review-final-489

\chapter{Cold Nights and Clear Skies}
\label{ch:ka-cold-nights-and-clear-skies}
\hlchaptermodality{489}

\begin{chapteropening}
\textbf{By the end of this chapter:} \emph{I can say dew, a chill wind, ice, to freeze and to twinkle.}

Complete the last lesson of chapter 489: I can say dew, a chill wind, ice, to freeze and to twinkle.
\end{chapteropening}

\section[\texorpdfstring{\kn{ಇಬ್ಬನಿ} (ibbani)}{ಇಬ್ಬನಿ (ibbani)}]{\kn{ಇಬ್ಬನಿ} (ibbani) — dew}

\label{lesson:KA-C489-ibbani}



\begin{quote}
Before the new one: say the Kannada for a crack, then the Kannada for to subside.
\end{quote}

\subsection*{You'll want to know: \kn{ಇಬ್ಬನಿ}}
\textbf{\kn{ಇಬ್ಬನಿ}} — \emph{ibbani} — \textquotedblleft{}dew\textquotedblright{}.

\begin{quote}
\kn{ಬೆಳಗ್ಗೆ} \kn{ಹುಲ್ಲಿನ} \kn{ಮೇಲೆ} \kn{ಇಬ್ಬನಿ} \kn{ಇತ್ತು}.
\end{quote}

\begin{quote}
\emph{beḷagge hullina mēle ibbani ittu.}
\end{quote}

\begin{quote}
In the morning there was dew on the grass.
\end{quote}

Drops on leaves at dawn after a cold night.

\begin{cousinweb}
\textbf{\kn{ಇಬ್ಬನಿ}} (\emph{ibbani}) is native Dravidian: DEDR 522 lists Kannada \emph{ibbani}, \textquotedblleft{}fog, mist, dew\textquotedblright{}.
\end{cousinweb}

\begin{grammarlens}[title={What you've built}]
One more word for getting around.
\end{grammarlens}

\subsection*{Your turn}
\begin{itemize}
\raggedright
  \item \emph{Say it:} \emph{ibbani}
  \item \emph{Say it:} \emph{ibbani}, once more
  \item \emph{Say it:} \emph{beḷagge hullina mēle ibbani ittu.} — In the morning there was dew on the grass.
  \item \emph{Recall:} say the Kannada for a crack, then the Kannada for to subside, then say \emph{ibbani} again
\end{itemize}

\begin{tcolorbox}[breakable,colback=teal!4,colframe=teal!35!black,title={Before you move on}]
What does \kn{ಇಬ್ಬನಿ} mean? (\textquotedblleft{}Dew\textquotedblright{}.) Say it once more.
\end{tcolorbox}
\section[\texorpdfstring{\kn{ಕುಳಿರ್ಗಾಳಿ} (kuḷirgāḷi)}{ಕುಳಿರ್ಗಾಳಿ (kuḷirgāḷi)}]{\kn{ಕುಳಿರ್ಗಾಳಿ} (kuḷirgāḷi) — a chill wind}

\label{lesson:KA-C489-kulirgali}



\begin{quote}
Before the new one: say the Kannada for to subside, then the Kannada for dew.
\end{quote}

\subsection*{You'll want to know: \kn{ಕುಳಿರ್ಗಾಳಿ}}
\textbf{\kn{ಕುಳಿರ್ಗಾಳಿ}} — \emph{kuḷirgāḷi} — \textquotedblleft{}a chill wind\textquotedblright{}.

\begin{quote}
\kn{ರಾತ್ರಿ} \kn{ಬೆಟ್ಟದ} \kn{ಮೇಲೆ} \kn{ಕುಳಿರ್ಗಾಳಿ} \kn{ಬೀಸುತ್ತಿತ್ತು}.
\end{quote}

\begin{quote}
\emph{rātri beṭṭada mēle kuḷirgāḷi bīsuttittu.}
\end{quote}

\begin{quote}
At night a chill wind was blowing on the hill.
\end{quote}

\kn{ಕುಳಿರ್}, cold, plus \kn{ಗಾಳಿ}, wind.

\begin{cousinweb}
\textbf{\kn{ಕುಳಿರ್ಗಾಳಿ}} (\emph{kuḷirgāḷi}) is a native Kannada compound of \emph{kuḷir}, \textquotedblleft{}to be cool or cold; coldness\textquotedblright{} (DEDR 1834), and \kn{ಗಾಳಿ}, \textquotedblleft{}wind\textquotedblright{} (DEDR 1499).
\end{cousinweb}

\begin{grammarlens}[title={What you've built}]
One more word for getting around.
\end{grammarlens}

\subsection*{Your turn}
\begin{itemize}
\raggedright
  \item \emph{Say it:} \emph{kuḷirgāḷi}
  \item \emph{Say it:} \emph{kuḷirgāḷi}, once more
  \item \emph{Say it:} \emph{rātri beṭṭada mēle kuḷirgāḷi bīsuttittu.} — At night a chill wind was blowing on the hill.
  \item \emph{Recall:} say the Kannada for to subside, then the Kannada for dew, then say \emph{kuḷirgāḷi} again
\end{itemize}

\begin{tcolorbox}[breakable,colback=teal!4,colframe=teal!35!black,title={Before you move on}]
What does \kn{ಕುಳಿರ್ಗಾಳಿ} mean? (\textquotedblleft{}A chill wind\textquotedblright{}.) Say it once more.
\end{tcolorbox}
\section[\texorpdfstring{\kn{ಮಂಜುಗಡ್ಡೆ} (mañjugaḍḍe)}{ಮಂಜುಗಡ್ಡೆ (mañjugaḍḍe)}]{\kn{ಮಂಜುಗಡ್ಡೆ} (mañjugaḍḍe) — ice; a block of ice}

\label{lesson:KA-C489-manjugadde}



\begin{quote}
Before the new one: say the Kannada for dew, then the Kannada for a chill wind.
\end{quote}

\subsection*{You'll want to know: \kn{ಮಂಜುಗಡ್ಡೆ}}
\textbf{\kn{ಮಂಜುಗಡ್ಡೆ}} — \emph{mañjugaḍḍe} — \textquotedblleft{}ice; a block of ice\textquotedblright{}.

\begin{quote}
\kn{ಚಳಿಗಾಲದಲ್ಲಿ} \kn{ಪರ್ವತದ} \kn{ಮೇಲೆ} \kn{ಮಂಜುಗಡ್ಡೆ} \kn{ಇರುತ್ತದೆ}.
\end{quote}

\begin{quote}
\emph{caḷigāladalli parvatada mēle mañjugaḍḍe iruttade.}
\end{quote}

\begin{quote}
In winter there is ice on the mountain.
\end{quote}

\kn{ಮಂಜು}, mist or frost, plus \kn{ಗಡ್ಡೆ}, a lump.

\begin{cousinweb}
\textbf{\kn{ಮಂಜುಗಡ್ಡೆ}} (\emph{mañjugaḍḍe}) is a native Kannada compound of \kn{ಮಂಜು}, \textquotedblleft{}dew, fog, coldness\textquotedblright{} (DEDR 4641), and \emph{gaḍḍe}, \textquotedblleft{}a mass, a lump\textquotedblright{} (DEDR 1148).
\end{cousinweb}

\begin{grammarlens}[title={What you've built}]
One more word for getting around.
\end{grammarlens}

\subsection*{Your turn}
\begin{itemize}
\raggedright
  \item \emph{Say it:} \emph{mañjugaḍḍe}
  \item \emph{Say it:} \emph{mañjugaḍḍe}, once more
  \item \emph{Say it:} \emph{caḷigāladalli parvatada mēle mañjugaḍḍe iruttade.} — In winter there is ice on the mountain.
  \item \emph{Recall:} say the Kannada for dew, then the Kannada for a chill wind, then say \emph{mañjugaḍḍe} again
\end{itemize}

\begin{tcolorbox}[breakable,colback=teal!4,colframe=teal!35!black,title={Before you move on}]
What does \kn{ಮಂಜುಗಡ್ಡೆ} mean? (\textquotedblleft{}Ice; a block of ice\textquotedblright{}.) Say it once more.
\end{tcolorbox}
\section[\texorpdfstring{\kn{ಹೆಪ್ಪುಗಟ್ಟು} (heppugaṭṭu)}{ಹೆಪ್ಪುಗಟ್ಟು (heppugaṭṭu)}]{\kn{ಹೆಪ್ಪುಗಟ್ಟು} (heppugaṭṭu) — to freeze; to set}

\label{lesson:KA-C489-heppugattu}



\begin{quote}
Before the new one: say the Kannada for a chill wind, then the Kannada for ice.
\end{quote}

\subsection*{You'll want to know: \kn{ಹೆಪ್ಪುಗಟ್ಟು}}
\textbf{\kn{ಹೆಪ್ಪುಗಟ್ಟು}} — \emph{heppugaṭṭu} — \textquotedblleft{}to freeze; to set\textquotedblright{}.

\begin{quote}
\kn{ರಾತ್ರಿ} \kn{ಚಳಿಗೆ} \kn{ಬಕೆಟ್ನ} \kn{ನೀರು} \kn{ಹೆಪ್ಪುಗಟ್ಟಿತು}.
\end{quote}

\begin{quote}
\emph{rātri caḷige bakeṭna nīru heppugaṭṭitu.}
\end{quote}

\begin{quote}
In the cold of the night the water in the bucket froze.
\end{quote}

Past (it …): \kn{ಹೆಪ್ಪುಗಟ್ಟಿತು}, \emph{heppugaṭṭitu}.

In the kitchen it is milk setting into curd: \kn{ಹೆಪ್ಪು} is the starter.

\begin{cousinweb}
\textbf{\kn{ಹೆಪ್ಪುಗಟ್ಟು}} (\emph{heppugaṭṭu}) is a native Kannada compound of \emph{heppu}, \textquotedblleft{}a curdling agent\textquotedblright{} (DEDR 4421), and \kn{ಕಟ್ಟು}, \textquotedblleft{}to gather, to harden\textquotedblright{} (DEDR 1148).
\end{cousinweb}

\begin{grammarlens}[title={What you've built}]
One more word for getting around.
\end{grammarlens}

\subsection*{Your turn}
\begin{itemize}
\raggedright
  \item \emph{Say it:} \emph{heppugaṭṭu}
  \item \emph{Say it:} \emph{heppugaṭṭu}, once more
  \item \emph{Say it:} \emph{rātri caḷige bakeṭna nīru heppugaṭṭitu.} — In the cold of the night the water in the bucket froze.
  \item \emph{Recall:} say the Kannada for a chill wind, then the Kannada for ice, then say \emph{heppugaṭṭu} again
\end{itemize}

\begin{tcolorbox}[breakable,colback=teal!4,colframe=teal!35!black,title={Before you move on}]
What does \kn{ಹೆಪ್ಪುಗಟ್ಟು} mean? (\textquotedblleft{}To freeze; to set\textquotedblright{}.) Say it once more.
\end{tcolorbox}
\section[\texorpdfstring{\kn{ಮಿನುಗು} (minugu)}{ಮಿನುಗು (minugu)}]{\kn{ಮಿನುಗು} (minugu) — to twinkle; to glitter}

\label{lesson:KA-C489-minugu}



\begin{quote}
Before the new one: say the Kannada for ice, then the Kannada for to freeze.
\end{quote}

\subsection*{You'll want to know: \kn{ಮಿನುಗು}}
\textbf{\kn{ಮಿನುಗು}} — \emph{minugu} — \textquotedblleft{}to twinkle; to glitter\textquotedblright{}.

\begin{quote}
\kn{ಕತ್ತಲೆಯಲ್ಲಿ} \kn{ಆಕಾಶದಲ್ಲಿ} \kn{ಚುಕ್ಕಿಗಳು} \kn{ಮಿನುಗುತ್ತಿದ್ದವು}.
\end{quote}

\begin{quote}
\emph{kattaleyalli ākāśadalli cukkigaḷu minuguttiddavu.}
\end{quote}

\begin{quote}
In the dark the stars were twinkling in the sky.
\end{quote}

Past (it …): \kn{ಮಿನುಗಿತು}, \emph{minugitu}.

Of stars, lights and jewels.

\begin{cousinweb}
\textbf{\kn{ಮಿನುಗು}} (\emph{minugu}) is native Dravidian: DEDR 4876 lists Kannada \emph{miṇuku}, \emph{minuku}, \emph{minugu}, \textquotedblleft{}to glitter, to sparkle\textquotedblright{}, the family of \kn{ಮಿಂಚು}, lightning.
\end{cousinweb}

\begin{grammarlens}[title={What you've built}]
One more word for getting around.
\end{grammarlens}

\subsection*{Your turn}
\begin{itemize}
\raggedright
  \item \emph{Say it:} \emph{minugu}
  \item \emph{Say it:} \emph{minugu}, once more
  \item \emph{Say it:} \emph{kattaleyalli ākāśadalli cukkigaḷu minuguttiddavu.} — In the dark the stars were twinkling in the sky.
  \item \emph{Recall:} say the Kannada for ice, then the Kannada for to freeze, then say \emph{minugu} again
\end{itemize}

\begin{tcolorbox}[breakable,colback=teal!4,colframe=teal!35!black,title={Before you move on}]
What does \kn{ಮಿನುಗು} mean? (\textquotedblleft{}To twinkle; to glitter\textquotedblright{}.) Say it once more.
\end{tcolorbox}
\section[(dialogue)]{Getting around, again}

\label{lesson:KA-R489-review-483-489}



\begin{quote}
Say the Kannada for a reservoir and a dam.
\end{quote}

\subsection*{Your turn}
Say these aloud:

\begin{itemize}
\raggedright
  \item a reservoir, a dam, backwater, groundwater, to evaporate
  \item temperature, humidity, climate, dampness, to worsen
  \item the monsoon, a dark cloud, drizzle, a downpour, to get soaked
  \item the late monsoon, hail, haze, to shroud, snowfall
  \item a cyclone, a whirlwind, to rage, a landslide, a tsunami
  \item a heatwave, scorching sun, to run dry, a crack, to subside
  \item dew, a chill wind, ice, to freeze, to twinkle
\end{itemize}

\begin{tcolorbox}[breakable,colback=teal!4,colframe=teal!35!black,title={Before you move on}]
A reservoir? (\textbf{\kn{ಜಲಾಶಯ}}, \emph{jalāśaya}.) A dam? (\textbf{\kn{ಅಣೆಕಟ್ಟು}}, \emph{aṇekaṭṭu}.) Backwater? (\textbf{\kn{ಹಿನ್ನೀರು}}, \emph{hinnīru}.) Groundwater? (\textbf{\kn{ಅಂತರ್ಜಲ}}, \emph{antarjala}.) To evaporate? (\textbf{\kn{ಆವಿಯಾಗು}}, \emph{āviyāgu}.) Temperature? (\textbf{\kn{ತಾಪಮಾನ}}, \emph{tāpamāna}.) Humidity? (\textbf{\kn{ಆರ್ದ್ರತೆ}}, \emph{ārdrate}.) Climate? (\textbf{\kn{ಹವಾಗುಣ}}, \emph{havāguṇa}.) Dampness? (\textbf{\kn{ತೇವ}}, \emph{tēva}.) To worsen? (\textbf{\kn{ಹದಗೆಡು}}, \emph{hadageḍu}.) The monsoon? (\textbf{\kn{ಮುಂಗಾರು}}, \emph{muṅgāru}.) A dark cloud? (\textbf{\kn{ಕಾರ್ಮೋಡ}}, \emph{kārmōḍa}.) Drizzle? (\textbf{\kn{ತುಂತುರು}}, \emph{tunturu}.) A downpour? (\textbf{\kn{ಜಡಿಮಳೆ}}, \emph{jaḍimaḷe}.) To get soaked? (\textbf{\kn{ತೊಯ್ಯು}}, \emph{toyyu}.) The late monsoon? (\textbf{\kn{ಹಿಂಗಾರು}}, \emph{hiṅgāru}.) Hail? (\textbf{\kn{ಆಲಿಕಲ್ಲು}}, \emph{ālikallu}.) Haze? (\textbf{\kn{ಮಬ್ಬು}}, \emph{mabbu}.) To shroud? (\textbf{\kn{ಮುಸುಕು}}, \emph{musuku}.) Snowfall? (\textbf{\kn{ಹಿಮಪಾತ}}, \emph{himapāta}.) A cyclone? (\textbf{\kn{ಚಂಡಮಾರುತ}}, \emph{caṇḍamāruta}.) A whirlwind? (\textbf{\kn{ಸುಂಟರಗಾಳಿ}}, \emph{suṇṭaragāḷi}.) To rage? (\textbf{\kn{ಅಬ್ಬರಿಸು}}, \emph{abbarisu}.) A landslide? (\textbf{\kn{ಭೂಕುಸಿತ}}, \emph{bhūkusita}.) A tsunami? (\textbf{\kn{ಸುನಾಮಿ}}, \emph{sunāmi}.) A heatwave? (\textbf{\kn{ಬಿಸಿಗಾಳಿ}}, \emph{bisigāḷi}.) Scorching sun? (\textbf{\kn{ಉರಿಬಿಸಿಲು}}, \emph{uribisilu}.) To run dry? (\textbf{\kn{ಬತ್ತು}}, \emph{battu}.) A crack? (\textbf{\kn{ಬಿರುಕು}}, \emph{biruku}.) To subside? (\textbf{\kn{ತಗ್ಗು}}, \emph{taggu}.) Dew? (\textbf{\kn{ಇಬ್ಬನಿ}}, \emph{ibbani}.) A chill wind? (\textbf{\kn{ಕುಳಿರ್ಗಾಳಿ}}, \emph{kuḷirgāḷi}.) Ice? (\textbf{\kn{ಮಂಜುಗಡ್ಡೆ}}, \emph{mañjugaḍḍe}.) To freeze? (\textbf{\kn{ಹೆಪ್ಪುಗಟ್ಟು}}, \emph{heppugaṭṭu}.) To twinkle? (\textbf{\kn{ಮಿನುಗು}}, \emph{minugu}.)
\end{tcolorbox}
\section[(dialogue)]{Cold Nights and Clear Skies — once more}

\label{lesson:KA-R489-review-final-489}



\begin{quote}
Say the Kannada for dew and a chill wind.
\end{quote}

\subsection*{Your turn}
\begin{itemize}
\raggedright
  \item \emph{Say it:} dew, a chill wind, ice, to freeze, to twinkle
\end{itemize}

\begin{tcolorbox}[breakable,colback=teal!4,colframe=teal!35!black,title={Before you move on}]
Dew? (\textbf{\kn{ಇಬ್ಬನಿ}}, \emph{ibbani}.) A chill wind? (\textbf{\kn{ಕುಳಿರ್ಗಾಳಿ}}, \emph{kuḷirgāḷi}.) Ice? (\textbf{\kn{ಮಂಜುಗಡ್ಡೆ}}, \emph{mañjugaḍḍe}.) To freeze? (\textbf{\kn{ಹೆಪ್ಪುಗಟ್ಟು}}, \emph{heppugaṭṭu}.) To twinkle? (\textbf{\kn{ಮಿನುಗು}}, \emph{minugu}.)
\end{tcolorbox}
